\documentclass[../main.tex]{subfiles}
\begin{document}
%====TIÊU ĐỀ TẬP==
\begin{table}[h]
\begin{adjustwidth}{-1cm}{}
\begin{tabular}{>{\centering\arraybackslash}p{6cm}|>{\raggedright\arraybackslash}p{12.5cm}}
\multirow{5}{6cm}
{\\[-40pt]
\flaregothic\fontsize{20pt}{20pt}\selectfont \centering
\color{eptype}HISTOIRE\color{black}\\
\freshbot\fontsize{72pt}{72pt}\selectfont
\color{epnum}17\\[6pt]
\cafeteria\fontsize{20pt}{20pt}\selectfont\color{epnum}(D-17)\color{black}
\\[18pt]
\color{black}
}
&
{\fotskip\fontsize{18pt}{18pt}\selectfont \ruby{倉}{そう}\ruby{庫}{こ}の\ruby{臨}{りん}\ruby{時}{じ}\ruby{隊}{たい}\ruby{長}{ちょう}}
\\[6pt]
& {\cafeteria\fontsize{18pt}{18pt}\selectfont A Shift Without a Roar} \\[4pt]
& {\vnmsans\fontsize{18pt}{18pt}\selectfont Ca trực không cần tiếng gầm} \\[8pt]
& {\caviar{\fontsize{14pt}{14pt}\selectfont Nhân vật: \emp{kson} \emp{kuon2} \emp{kaigou2} \emp{game} \emp{delu} \emp{ryuuhou2} \emp{suzuno2} \emp{naomi2}}} \\[8pt]
\end{tabular}
\end{adjustwidth}
\end{table}
%===TRUYỆN CHÍNH==
\color{black}

\txtbx[2]{
{\color{vol} \uline{Tóm tắt:}} Khi kho vật liệu xây dựng của bà Kaien thiếu người vào một buổi sáng mưa lớn, Kson xung phong nhận ca phụ giúp. Một vệt nước mưa làm nhòe mã quy cách trên phiếu giao hàng khiến nhiều bó thép hộp, tấm tôn lạnh và hệ xương cá kim loại có nguy cơ bị chở nhầm xe — trong đó có đơn vật liệu mà một công trình đang chờ gấp để kịp lợp mái trước khi dầm mưa. Kson phối hợp cùng Kuon, Game và các tài xế xe nâng để rà soát kích thước, đối chiếu bản vẽ và điều phối lại tuyến xe tải giao hàng.
}

\pov{Kson}

\lang{Etsunan}

Buổi sáng hôm ấy bắt đầu bằng tiếng mưa lộp bộp gõ dồn dập trên mái tôn hiên nhà và một cuộc gọi từ bà Kaien.

Tôi đang ngồi ở bàn ăn, hai tay ôm cốc cà phê nóng nghi ngút khói, thì Kuon nhận điện thoại thay mẹ. Anh ấy chăm chú lắng nghe, thỉnh thoảng đáp ``vâng'' hoặc ``để con xem'', vẻ mặt bình thản đến mức tôi chưa đoán được chuyện gì.

Cho đến khi anh ấy cúp máy, xoay người lại nghiêng đầu nhìn tôi: ``Kho vật liệu xây dựng bên mẹ anh thiếu một người phụ đối chiếu mã hàng sáng nay. Mấy chuyến xe tải cẩu chở tôn và sắt hộp tới sớm hơn dự kiến, mà một anh phụ kho lại dở chứng bị sốt đột xuất. Em có rảnh không?''

Tôi đặt cạch cốc cà phê xuống bàn, hào hứng đập tay lên ngực: ``Rảnh chứ! Mấy việc khuân vác hay sắt thép nặng nhọc cứ giao hết cho em!''

Kuon nhìn tôi thêm một hồi, đôi mắt hiện lên vẻ đắn đo như thể đang cân nhắc xem tôi có hiểu hết quy mô của một kho vật liệu hay chưa: ``Không phải chỉ có vác nặng đâu. Kiểm đếm mã quy cách, đo độ dài thanh xà gồ, dán nhãn tuyến giao thôi. Mẹ anh chỉ cần thêm một người phụ tay phụ chân đối chiếu đơn cho chính xác.''

``Anh nói vậy nghe như em định đấm thủng mấy xấp tôn lạnh của mẹ anh không bằng ấy,'' tôi bĩu môi hờn dỗi.

``Anh chỉ đang tính đến mọi khả năng có thể xảy ra thôi,'' Kuon điềm tĩnh đáp.

Tôi bật cười, đứng dậy lấy chiếc áo khoác da treo gần cửa: ``Vậy thì yên tâm đi, hôm nay em sẽ cố không làm gãy thanh xương cá nào đâu!''

Kuon nhướng mày, ánh mắt tràn đầy sự hoài nghi: ``Chỉ ``cố'' thôi à?''

``Thì em không dám hứa quá khả năng của mình chứ sao!'' tôi nháy mắt tinh nghịch.

\ct

Kho vật liệu xây dựng Kaien nằm ở cuối một vùng đất rộng ngoại ô thành phố, cách nhà Hikawa tầm ba mươi cây số. Nơi đó rộng tút tắp với bãi chứa gạch thẻ và sắt thép ngoài trời có mái che cao vút. Lúc nào nơi đây cũng nhộn nhịp tiếng bíp bíp rào rào của xe nâng di chuyển, tiếng kim loại va chạm lách cách và mùi thép mạ kẽm hòa lẫn với mùi bụi xi măng khô đặc trưng.

Bà Kaien đón chúng tôi ngay trước cổng kho, tươi cười đưa cho tôi một chiếc áo khoác phản quang màu cam cùng đôi găng tay bảo hộ dính chút vôi xám rất vừa vặn: ``Cảm ơn cháu đến giúp nhé, Kson-chan. Hôm nay có mấy chuyến xe tải cần chở tôn lạnh và sắt hộp đi trước mười một giờ quá.''

``Cô cứ giao việc cho cháu! Cháu làm việc ở mấy bãi vật liệu nặng thế này quen tay lắm ạ,'' tôi xắn tay áo đầy tự tin.

Kuon đứng phía sau khẽ hắng giọng, định cất lời phân tích về các quy chuẩn an toàn lao động thì tôi đã quay lại giơ ngón tay cái cắt ngang. ``Em từng phụ kho cơ khí với bãi vật liệu nhiều lần rồi nhé,'' tôi nói thêm, liếc sang anh ấy với vẻ tự hào. ``Đừng có nhìn em như thể em vừa nhận nhiệm vụ đi vận hành cần cẩu tháp chứ.''

``Anh đã kịp nói gì đâu nào,'' Kuon thở dài, khóe môi khẽ giật giật. ``Mà nhắc tới bãi cơ khí, cái vụ em với Kana-chan nghịch ngợm hồi đó làm cho cả văn phòng nổ tung\footnote{{\flaregothic ÉPISODE 62}, quyển số Năm thuộc \textit{holo no Graffiti}.} anh vẫn chưa quên đâu đấy.''

``Ơ kìa! Đừng có mà nhắc lại vụ đó ngay trước mặt phụ huynh chứ\~{}'' tôi nhăn mặt phân trần.

Bà Kaien bật cười vui vẻ trước màn đối đáp của hai đứa, vẫy tay dẫn chúng tôi vào sâu bên trong khu nhà kho chính.

Ở trung tâm kho là những dãy kệ thép chịu lực hạng nặng chất đầy các cuộn tôn mạ màu, từng bó sắt hộp 40x80 bó đai thép gọn gàng và những xấp thanh xương cá làm trần thạch cao dài ngoẵng. Mỗi dãy được chia rõ ràng theo từng tuyến xe tải giao hàng. Những tấm bảng màu xanh, vàng và trắng treo phía trên giúp tài xế xe nâng dễ dàng nhận biết khu vực tập kết.

Ở góc sâu nhất của kho, tôi chú ý thấy một lô thép tấm gia cố chịu lực dày dặn cùng những kiện vật liệu cách âm đặc biệt được bọc bạt kín đáo, dán nhãn ghi chú riêng.

Thấy ánh mắt tôi nheo lại tò mò, Kuon khẽ giải thích: ``Lô thép chịu lực và hệ khung chuyên dụng đó là mẹ tích trữ riêng cho dự án gia cố công trình sắp tới của gia đình đấy. Đừng chạm vào kẻo lệch mã kiểm kê.''

Từ chiếc đồng hồ đeo tay của Kuon, giọng Game vang lên vừa đủ nghe, mang chút vẻ thong thả: ``Tôi đồng bộ xong toàn bộ bảng kê quy cách vật liệu sang máy quét cầm tay cho mọi người rồi đấy. Cứ làm đúng theo mã lô và độ dày tôn là được, đừng có tự ý sáng tạo kích thước sắt rồi lại bắt tôi sửa hệ thống.''

``Còn nếu mọi người làm theo lời tôi thì bảo đảm ca làm việc sẽ sôi động hơn nhiều!'' tôi toe tút cười lớn.

``Quan trọng là phải chính xác từng milimet độ dày, Kson-san, chứ không phải là mức độ sôi động của cô đâu,'' Game đáp trả không chút nể hờn.

Kuon đưa cho tôi một tập phiếu giao hàng cùng chiếc máy quét mã vạch: ``Em kiểm các đơn sắt hộp và tôn lợp đi tuyến phía Đông trước đi. Cứ quét mã đai thép, đối chiếu số bó và đánh dấu vào bảng kê. Đừng ra hiệu cho xe nâng bốc gì lên thùng xe tải cho đến khi người phụ trách xác nhận đấy.''

``Rõ, đội trưởng!'' tôi giơ tay chào theo phong cách quân đội.

``Anh không phải đội trưởng đâu,'' Kuon lờ đi trò đùa của tôi.

``Vậy ai đang ra lệnh cho em nào?''

Anh ấy thản nhiên chỉ tay về phía mẹ mình: ``Người trả lương cho em đấy.''

``Thế thì cô Kaien mới đúng là đội trưởng!'' Tôi nghiêm túc xoay người sang bà Kaien. ``Đội trưởng, cháu đã sẵn sàng nhận lệnh!''

Bà Kaien mỉm cười xoa đầu tôi, vẫy tay chỉ về phía bãi tập kết tôn lạnh: ``Vậy đội phó Kson bắt đầu kiểm tra tuyến giao cho công trình An Wa trước nhé.''

\ct

Trong nửa giờ đầu tiên, mọi thứ diễn ra êm đẹp đến bất ngờ.

Tôi quét mã từng bó sắt, đọc to kích thước và độ dày tôn trên phiếu rồi nhịp nhàng ra hiệu cho chiếc xe nâng màu vàng từ từ hạ càng nâng, di chuyển kiện hàng vào đúng vị trí bãi chờ xe tải. Những xấp thanh xương cá trần nhẹ hơn thì tôi cùng một anh công nhân tự tay khiêng xếp gọn gàng sang bên cạnh.

Tôi thao tác nhanh gọn nhưng hoàn toàn không hấp tấp. Một lần Kuon tiến lại gần định nhắc tôi kiểm tra lại quy cách xà gồ C, tôi đã nhanh tay giơ thước kẹp cơ khí cùng tờ phiếu lên trước mặt anh ấy trước khi anh ấy kịp mở miệng.

``Đơn VL-03, tổng cộng bốn bó sắt hộp mạ kẽm và hai cuộn tôn lạnh 9 sóng. Đúng độ dày 0.45mm và chuẩn mã quy cách!'' tôi dõng dạc báo cáo.

``Tốt lắm,'' Kuon gật đầu.

``Anh nói câu đó nghe như thể đang ngạc nhiên lắm ấy nhỉ?'' tôi nhướng mày trêu chọc.

``Anh chỉ đang cập nhật lại đánh giá về em thôi. Chuyện bình thường mà,'' Kuon thản nhiên đáp.

``Cập nhật từ ``thích phá hoại'' lên ``biết phân biệt sắt hộp với tôn lạnh'' đúng không?'' tôi cười khì khì.

Kuon nhún vai từ tốn: ``Đại loại thế.''

Tôi bật cười sảng khoái, xoay người tiếp tục làm việc. Đúng lúc đó, tiếng động cơ ầm ầm vang lên, chiếc xe tải thùng phủ bạt thứ hai từ từ lùi vào khoảng sân B trước kho.

Cơn mưa ngoài bãi lúc này đã bắt đầu nặng hạt. Bác tài xế bước xuống cabin, vai áo mưa đã ướt đẫm, hai tay ôm xấp giấy phiếu giao hàng được bọc vội trong túi nhựa nhưng vẫn bị hắt nước: ``Xin lỗi cô Kaien, đoạn đường tránh phía Đông bị ngập sâu nên xe tới trễ. Mấy cái nhãn dán mã quy cách trên các bó tôn bị nước mưa tạt trôi nhòe hết mực rồi!''

Kuon nhanh bước ra nhận lấy bảng kê. Mẹ anh nhìn lên chiếc đồng hồ treo tường rồi thở dài lo lắng: ``Nếu mất thêm thời gian đo đạc phân loại lại từng thanh sắt từ đầu thì công trình nhà xưởng An Wa sẽ không kịp có tôn lợp mái trước khi cơn mưa rào buổi trưa đổ xuống mất.''

Tôi đặt ngay máy quét xuống bàn, tiến lại gần: ``Cô cho cháu xem thử xấp phiếu được không ạ?''

Trên mặt giấy, mực bút lông đã loang lổ thành những vệt màu xanh tím. Một vài mã tuyến vẫn còn đọc được, nhưng phần quy cách chiều dài (4 mét hay 6 mét) của ba lô vật liệu quan trọng nhất gần như đã biến thành những đám mây nhòe nhoẹt.

Chiếc đồng hồ trên tay Kuon phát ra tiếng *tít* nhẹ. Game tặc lưỡi đánh giá: ``Hai mã đầu thì tôi tra ngược lại lịch sử xuất kho được, nhưng lô thứ ba bị vệt nước làm nhòe tơi tả rồi. Đừng có nhìn tôi bằng ánh mắt kỳ vọng thế, tôi là trợ lý thông minh chứ có phải kính hiển vi soi vết mực tan đâu.''

``Vậy thì đừng đoán mò,'' Kuon lập tức đưa ra quyết định. ``Chúng ta sẽ xác nhận bằng phương pháp đối chiếu trực tiếp.''

Tôi gật đầu đồng tình, giơ điện thoại lên: ``Để em gọi trực tiếp cho cai thầu công trình An Wa đối chiếu số lượng và bản vẽ cho chắc!''

\ct

Bà Kaien nhanh chóng lấy sổ theo dõi đơn hàng bản giấy từ văn phòng kho. Kuon mở danh sách đối chiếu trên máy tính bảng, còn tôi chịu trách nhiệm quay số cho bên thi công An Wa. Bên kia đầu dây, tiếng máy cắt sắt và tiếng búa đập chát chúa vang lên dồn dập: ``Alo! Đội thi công An Wa nghe đây ạ!''

``Dạ cháu chào chú, cháu là Kson gọi từ kho vật liệu xây dựng Kaien đây ạ!'' Tôi dồn hết sự nhiệt tình vào giọng nói. ``Cháu đang kiểm tra lại chuyến xe giao tôn và sắt sáng nay. Đội mình đang chờ tôn lạnh 9 sóng, sắt hộp 40x80 với hệ xương cá trần đúng không chú?''

``Đúng rồi cháu ơi! Hai mươi tấm tôn dài 6 mét, bốn bó sắt hộp mạ kẽm với năm xấp xương cá. Mưa sắp kéo tới to rồi, xe tải chở tôn tới nơi chưa cháu? Mấy thợ đang đứng trên giàn giáo chờ lợp mái đây!''

``Bọn cháu đang cho xe nâng xếp lại đúng thùng xe cho chuẩn xác. Cháu sẽ báo bác tài xuất bến ngay trong năm phút nữa nhé!''

Tôi cúp máy, dùng bút đánh dấu dầu ghi lại chính xác thông số ``Tôn 6m - 20 tấm'' lên tấm bìa cứng rồi dán chặt vào đai thép của bó tôn.

Đơn thứ hai thuộc về một xưởng cơ khí nhỏ gần bến xe, họ xác nhận đặt thép góc V và hai mươi bao xi măng. Đơn thứ ba là của cửa hàng đại lý vật liệu, chủ tiệm xác nhận họ chờ một kiện đinh vít tự khoan và xà gồ C.

Kuon tỉ mỉ đối chiếu từng con số tôi vừa ghi với dữ liệu trong sổ theo dõi, rồi kết luận: ``Vậy ba lô bị nhòe là đơn tôn dài của công trình An Wa và xà gồ của đại lý.''

``Thế còn bó sắt hộp to đùng nằm ở góc xe nâng này thì sao?'' tôi vỗ vỗ tay lên một bó sắt mạ kẽm sáng bóng.

Game lập tức quét mã dán trên đai kẹp sắt rồi cất giọng: ``Đơn đó của công trình An Wa đấy. Sắt hộp 40x80 dày 1.8mm, đúng bốn bó như cai thầu vừa đọc.''

Tôi nheo mắt nhìn lại mảnh giấy ghi chép: ``Nhưng bên An Wa đặt bốn bó sắt, hai mươi tấm tôn và năm xấp xương cá. Mấy xấp xương cá trần đâu rồi nhỉ?''

Bà Kaien vội vã đi sang khu vực kệ cao. Quả nhiên, ba xấp xương cá kim loại đang nằm chèn phía sau mấy kiện gạch của đại lý. Hóa ra trong lúc tài xế xe nâng bốc xếp dưới mưa, nhãn dán dính nước bị bong ra, khiến nhân viên tiện tay xếp chúng cạnh nhau theo thứ tự phân loại thô chứ không theo đơn công trình.

``Nếu mình chỉ nhìn lướt qua thì tí nữa xe tải đã chở nhầm xương cá trần sang bãi gạch rồi,'' tôi tặc lưỡi lắc đầu.

Kuon cẩn thận dùng thước dây đo lại chiều dài thanh sắt rồi gật gù: ``Ừ, may mà em phát hiện kiểm lại. Giờ mình ra hiệu cho xe nâng chuyển riêng các lô ra rồi dán lại thẻ tuyến.''

``Game này, ông phát lệnh in lại mấy cái thẻ mã vạch mới cho tôi được không?'' tôi quay sang nhìn mặt đồng hồ.

Game đáp lại bằng giọng tỉnh rụi: ``Mẫu thẻ thì tôi chuyển sang máy in ở văn phòng kho rồi. Nhưng nhắc trước nhé, tôi không có tay để cầm kềm bấm đai dán hộ đâu đấy, đừng có ngồi đó mà trông chờ.''

``Tôi có tay mà, không mượn đến tay ông đâu nhé!'' tôi toe tút đáp trả.

\ct

Cả đội nhanh chóng dồn các kiện vật liệu chưa rõ nhãn vào khu vực khô ráo rồi bắt đầu quy trình bốc dỡ điều phối lại.

Tôi đọc to quy cách tôn thép, Kuon nhanh tay đối chiếu mã số trên hệ thống, Game phát tiếng bíp xác nhận số lượng, còn bà Kaien đánh dấu chéo vào sổ giao dịch. Mỗi khi cả bốn nguồn dữ liệu khớp nhau, tôi liền bấm đai thẻ nhựa mới tinh lên bó sắt rồi ra hiệu cho xe nâng cẩu nó sang thùng xe tải.

Làm việc tỉ mỉ với sắt thép nặng nề tuy tốn sức hơn việc ước lượng bằng mắt, nhưng đổi lại không một chuyến xe nào bị chở sai chủng loại.

Đúng lúc xe nâng xếp xong kiện cuối cùng thì điện thoại Kuon reo lên. Tiếng Kaigou vang lên rõ ràng qua loa ngoài: ``Aniki, bên kho vật liệu mọi người ổn cả chứ? Mẹ nhắn tin bảo mưa làm nhòe hết phiếu giao tôn thép rồi.''

``Đang xử lý xong rồi em. Có Kson-chan ở đây phụ điều phối xe nâng và đối chiếu mã tôn nên rất nhanh,'' Kuon trả lời.

Tôi vội ghé sát vào điện thoại, vẫy vẫy tay dù biết Kaigou không thấy được: ``Kaigou-paisen! Mọi thứ vẫn hoàn toàn trong tầm kiểm soát nhé! Mấy bó sắt hộp nặng trịch em cân đẹp hết!''

``Em chắc không đấy?'' giọng Kaigou đầy vẻ nghi vấn.

``Chắc chắn một trăm phần trăm! Bọn em đã gọi điện kiểm tra kích thước với bên thi công từng đơn một rồi!''

``Thế thì tốt. Chị đang ở gần khu đó, cần kéo xe hay bốc vác hàng nặng thì gọi chị. Nhưng mà Aniki đừng có để Kson tự tay vác cả cuộn tôn nặng nửa tấn đấy nhé. Trông to con thế thôi chứ thể lực không phải cái gì cũng nhấc nổi đâu.''

Tôi nhíu mày, hét lên vào điện thoại đầy vẻ bất bình: ``Em nghe thấy hết đấy nhé, Kaigou-paisen!''

``Chị cố ý nói cho em nghe mà,'' chị Tổng cười khúc khích rồi cúp máy.

Ngay sau đó, Ryuuhou gửi vào nhóm chat gia đình tấm ảnh chụp lại màn hình cập nhật tiến độ bốc dỡ vật liệu của Game, kèm dòng tin nhắn trêu đùa: \textit{``Game làm quản lý dữ liệu từ xa, Kson-san làm chỉ huy bãi xe nâng. Phân công lao động hợp lý phết!''}

Tôi nhanh tay gửi lại sticker hình ngón tay cái hào hứng.

Kuon liếc nhìn màn hình điện thoại của tôi: ``Em có thể nhắn là cả hai chúng ta cùng làm mà.''

``Nhắn như em thì nghe mới ngầu chứ!'' tôi nháy mắt.

\ct

Mười giờ hai mươi phút, chuyến xe tải thùng đầu tiên chở đầy tôn lạnh và sắt hộp đã sẵn sàng nổ máy rời kho.

Công trình An Wa nhận đủ tôn lợp mái và sắt khung; xưởng cơ khí bến xe nhận chuẩn xác thép V và xi măng; các đại lý cũng có đủ đinh vít và xà gồ. Tuy nhiên, chiếc xe tải nhỏ chuyên chở hệ xương cá nhẹ lại chưa thể quay về kho vì tuyến đường ngắn qua cầu sắt bị ngập sâu, buộc xe phải chạy đường vòng rất xa.

Bà Kaien lo lắng nhìn bảng giờ thi công của khách, khẽ mím môi: ``Nếu đi đường vòng thì e là đội thợ An Wa không kịp có xương cá để làm trần trước khi tạnh mưa mất\dots''

``Họ cần hệ xương cá đó gấp lắm ạ?'' tôi hỏi.

``Mười một giờ là họ bắt đầu bắn khung trần rồi. Mấy thanh kim loại này phải tới trước đó để thợ cắt quy cách.''

Tôi lập tức tiến lại gần bản đồ quy hoạch khu vực dán trên tường văn phòng kho: ``Từ đây đến công trình An Wa không xa lắm. Nếu xe tải lớn không đi được đường cầu ngập, mình có thể lách qua con đường bê tông nhỏ chạy song song với đê mà!''

Kuon tiến lại nhìn theo ngón tay tôi chỉ rồi chau mày: ``Đường đê đó hẹp lắm, xe tải thùng lớn chở tôn dài không vào nổi đâu.''

``Thế thì mình chuyển riêng mấy xấp thanh xương cá nhẹ sang chiếc xe bán tải chở hàng của kho!'' tôi đề xuất.

Bà Kaien gật đầu tán thành: ``Xe bán tải đang đậu ở bãi sau. Nhưng tài xế lại đang lái chiếc xe cẩu rồi.''

``Cháu đi cùng cô! Cháu thuộc đường đê và có thể phụ bê mấy thanh xương cá xuống công trình!'' tôi xung phong ngay.

Kuon lập tức chớp mắt nhìn tôi: ``Ê, này đừng tự ý đổi tuyến. Phải gọi xác nhận với tài xế xe tải trước đã.''

``Em đâu có nói là em tự lái xe bán tải lao đi mất đâu nào!'' tôi phân trần.

``Anh nhắc trước thế thôi, em làm gì đã có bằng lái xe tải hạng C!''

Mặt đồng hồ Kuon lại sáng lên, giọng Game cất lên cắt ngang: ``Tôi vừa kết nối với tài xế xe tải rồi. Anh ấy bảo có thể hội quân với xe bán tải ở ngã ba đường đê phía Nam sau khi trút xong đơn sắt hộp đầu tiên.''

Bà Kaien lập tức đưa ra quyết định: ``Vậy Kson đi cùng cô nhé! Chúng ta chở riêng hệ xương cá trần, còn Kuon, con ở lại kiểm tra nốt phần thép tấm dán nhãn lưu kho.''

``Rõ, đội trưởng!'' / ``Vâng ạ!''

\ct

Chiếc xe bán tải chở hàng màu xám xuất phát từ bãi sau, thùng xe phía sau xếp gọn năm xấp thanh xương cá kim loại sáng loáng. Bà Kaien cầm lái, còn tôi ngồi ở ghế phụ, tay ôm tập bản vẽ quy cách và rà soát lại số lượng lần cuối.

Con đường bê tông ven đê sau trận mưa sáng trở nên sạch sẽ lạ thường, mặt đường ẩm ướt phản chiếu ánh nắng vừa hửng lên qua kẽ mây.

Tại ngã ba, chiếc xe tải của kho đã chờ sẵn. Bác tài hạ kính xe, cười tươi đưa cho bà ấy tập phiếu đối chiếu đã khô ráo: ``Tôi xong tuyến sắt hộp rồi nhé! Đơn xương cá trần trông cậy vào hai cô cháu đấy!''

``Cảm ơn anh, bọn tôi đi ngay đây!''

Tôi chỉ tay về phía con đường ven đê: ``Cô đi hướng này cho nhanh ạ!''

Con đường hẹp đúng như Kuon nói, nhưng chiếc xe bán tải di chuyển rất dễ dàng. Hai bên đường, các cửa hàng đại lý và công trình xây dựng nhỏ đang tưng bừng mở cửa đón nắng mới, một xưởng mộc bày những tấm gỗ thơm phức ra hiên, còn quán mì đầu phố đã treo bảng hiệu nghi ngút khói.

Tôi nhận ra vài gương mặt quen thuộc trong khu phố. Một chú lái xe ba gác chở gạch vẫy tay chào khi thấy xe bán tải của kho Kaien chạy qua.

``Cháu Kson lại đi giao vật liệu đấy à?!'' ông chú gọi lớn.

``Hôm nay cháu làm phụ tá kho vật liệu xây dựng của cô Kaien ạ!'' tôi nhoài người ra cửa kính vẫy tay đáp lại rạng rỡ.

Bà Kaien bật cười vui vẻ: ``Cháu quen nhiều người quanh khu bãi này thật đấy Kson.''

``Cháu hay đi lượn mua đồ ăn với ngó nghiêng mấy xưởng cơ khí quanh khu này mà cô. Với lại mặt cháu cũng dễ nhớ nữa!''

``Có lẽ vì tiếng chào của cháu lúc nào cũng vang xa nhất phố đấy,'' bà trêu.

``Cô cũng nghe thấy tiếng cháu chào ạ?'' tôi ngượng ngùng gãi đầu.

``Cả cái khu bãi vật liệu này nghe thấy thì có!'' bà Kaien cười lớn.

\ct

Công trình nhà xưởng An Wa nằm ở góc đường mới mở, nổi bật với hàng giàn giáo thép dựng đứng và mái tôn vừa lợp xong một nửa. Bác cai thầu mặc chiếc áo bảo hộ hộ dính bẩn đang đứng ngồi không yên dưới hiên, mắt liên tục liếc nhìn đồng hồ.

``May quá! Xương cá trần tới rồi!'' bác reo lên khi thấy chiếc xe bán tải dừng lại.

``Cháu xin lỗi chú vì sự cố nhòe phiếu sáng nay nhé! Bọn cháu đã kiểm tra kỹ độ dài từng thanh một rồi ạ!''

Tôi nhanh nhẹn nhảy xuống xe, mở thùng sau rồi cẩn thận cùng bà Kaien nhấc xấp thanh xương cá kim loại đầu tiên xuống. Bác cai thầu nhanh chóng đối chiếu tấm thẻ mới dán trên từng xấp rồi gật đầu lia lịa: ``Đủ và chuẩn quy cách hết rồi! Tôi cứ lo sáng nay thợ phải ngồi chơi xào nước vì thiếu khung trần chứ!''

``Bọn cháu đã hứa là phải kịp mà chú!'' tôi mỉm cười tự hào.

Bác cai thầu nhìn tôi một hồi, rồi bỗng ngập ngừng hỏi: ``Có phải cháu là cô bé hay chào rõ to mỗi lần ghé mua nước ở đại lý vật liệu đầu phố không?''

Tôi khựng lại một chút, gãi gãi má: ``Dạ\dots\ là cháu ạ.''

``Tôi nhận ra cái giọng sảng khoái này ngay! Sáng nay nghe tiếng cháu gọi điện xác nhận kích thước sắt, tôi cũng thấy yên tâm hẳn,'' bác mỉm cười dịu dàng.

``Vậy lần sau cháu sẽ cố nói nhỏ lại một chút để chú khỏi phải giật mình nhé,'' tôi cười khì khì.

``Không cần đâu! Công trình này ồn ào tiếng máy cắt lắm, tiếng chào rõ ràng của cháu nghe phấn chấn tinh thần thợ lắm đấy.''

Bác cai thầu dẫn chúng tôi vào khu vực văn phòng điều hành tạm thời của công trình. Bà chủ xưởng ở đó nhanh tay rót hai ly trà đá mát lạnh, rồi lấy ra một đĩa bánh ngọt cùng vài trái cây tươi dúi vào tay tôi: ``Coi như quà cảm ơn người giao vật liệu đặc biệt đúng giờ nhé!''

``Ôi, cháu không thể nhận quà của khách hàng đâu cô chú ơi!'' tôi xua tay từ chối.

``Đây không phải hàng hóa mua bán, đây là chút quà vắt vai cô tặng riêng cho cháu vì đã nhiệt tình cứu tiến độ lợp mái,'' bà chủ mỉm cười kiên quyết.

Bà Kaien nhìn tôi, gật đầu khẽ: ``Cháu nhận đi Kson, họ đã nghiệm thu đủ đơn tôn thép rồi mà.''

Tôi hai tay nhận lấy dĩa bánh ấm áp, cúi đầu thật thấp: ``Vậy cháu cảm ơn cô chú nhiều ạ!''

\ct

Khi chúng tôi trở lại kho vật liệu, chiếc xe tải lớn cũng vừa hoàn thành xong lộ trình còn lại. Không còn lô tôn thép nào bị nhầm lẫn, không có công trình nào phải chậm tiến độ, và xấp phiếu giao hàng mới đã được kẹp ngay ngắn vào bìa hồ sơ chống nước.

Kuon đang ngồi trước máy tính trong văn phòng kho, gõ bàn phím liên tục để cập nhật trạng thái tồn kho sắt thép. Nghe tiếng bước chân, cậu ngẩng đầu lên: ``Giao xong hết rồi à em?''

``Xong xuôi hết rồi nhé! Công trình An Wa nhận đủ tôn, sắt hộp với xương cá trần trước mười một giờ luôn!'' tôi tự hào tuyên bố.

``Khá đấy. Tưởng em chỉ biết ra hiệu cho xe nâng chạy lung tung thôi chứ,'' Kuon trêu ghẹo.

Tôi đặt đĩa bánh cùng túi trái cây lên bàn làm việc của cậu: ``Và em được người ta trả công bằng cái này đây!''

Kuon nhìn đĩa bánh, nhướng mày nghi vấn: ``Anh tưởng đơn hàng đã được kiểm đủ rồi chứ? Em lại nhặt trộm vật liệu gì của công trình đổi lấy bánh đấy à?''

``Đây là quà chủ xưởng tặng riêng cho em, không phải hàng đổi hàng đâu nhé!'' tôi bĩu môi thanh minh.

Mặt đồng hồ của Kuon lại nhấp nháy sáng. Game lên tiếng bênh vực đầy chí lý: ``Tôi xác nhận đĩa bánh này hoàn toàn không có trong danh mục kiểm kê tôn thép. Nhưng nếu Kson-san có ý định giấu bớt một quả táo để ăn mảnh thì tôi cũng không ngại báo cáo lại với bà Kaien đâu nhé.''

``Game! Ông đúng là cái đồ cơ hội!'' tôi gào lên vui vẻ.

Bà Kaien bước ra từ khu vực kho trong, mím môi mỉm cười nhìn ba chúng tôi: ``Hôm nay nhờ có Kson phụ kiểm đếm mà chuyến giao vật liệu trôi chảy hơn cô tưởng nhiều đấy. Nhất là cái khoản đọc vanh vách độ dày tôn với quy cách sắt hộp nữa.''

``Cháu chỉ gọi điện với bấm đai lại mấy cái thẻ thôi ạ! Lần sau kho có thiếu người phụ xe nâng hay đối chiếu thép cứ gọi cháu nhé!''

Kuon liếc nhìn về phía góc kho nơi cất giữ lô thép chịu lực gia cố đặc biệt, khẽ mỉm cười: ``Yên tâm, sắp tới gia đình mình có dự án thi công lớn, lúc đó chắc chắn sẽ cần đến tay nghề kiểm đếm vật liệu của em đấy.''

Bà Kaien gật đầu hiền hậu: ``Được rồi, đến lúc đó cô sẽ bảo bên thi công chuẩn bị phần bánh gấp đôi cho cháu nhé!''

%====TRUYỆN PHỤ===

\omk

Buổi chiều muộn, tôi thong thả đi bộ về nhà Hikawa thay vì đi xe buýt.

Cơn mưa buổi sáng đã tạnh hẳn, nhường chỗ cho những vệt nắng vàng ươm vắt qua các mái nhà. Mặt đường còn ẩm ướt phản chiếu bầu trời trong trẻo. Tôi đi chậm qua xưởng mộc, đại lý vật liệu và cửa hàng tiện lợi quen thuộc. Ông chú lái xe ba gác lúc sáng lại giơ chai trà lạnh lên chào, tôi cũng vui vẻ vẫy tay đáp lại.

Chẳng có điều gì quá phi thường xảy ra trong ngày hôm nay cả. Tôi không giải cứu thế giới, cũng chẳng làm điều gì to tát ngoài việc xử lý vài tờ phiếu giao vật liệu bị nhòe mực và chỉ dẫn xe nâng bốc mấy bó sắt hộp. Nhưng cả ca làm việc trải qua với một nhịp điệu vô cùng dễ chịu và chắc chắn.

Khi tôi đẩy cửa bước vào nhà Hikawa, Delutaya đang ngồi trên sofa phòng khách, tay tháo một bên tai nghe ra nhìn tôi: ``Kson-san, chị mới đi đâu về thế ạ?''

``Chị vừa đi phụ kho vật liệu xây dựng về! Một ca điều phối tôn thép ngắn buổi sáng thôi,'' tôi vươn người giãn cơ.

``Làm ở kho vật liệu có vui không chị?'' Delutaya chớp mắt tò mò.

Tôi suy nghĩ một chút rồi cười tươi: ``Vui theo kiểu mệt nhoài cả tay chân, áo quần thì dính tí bụi xi măng, nhưng đổi lại được nghe tiếng xe nâng giòn giã và ăn bánh ngọt của chủ xưởng tặng!''

Delutaya bật cười khe khẽ: ``Nghe cũng thú vị đấy chứ.''

``Lần sau muốn đi xem xe nâng bốc tôn thì bảo chị, chị dẫn đi cùng. Nhưng phải đeo găng tay bảo hộ đấy nhé!''

Từ trên cầu thang, Ryuuhou chạy xồng xộc xuống, mũi hít hà: ``Kson-san! Chị đi kho về có mang bánh ngọt hay quà gì về cho em không đấy?!''

Tôi giơ chiếc đĩa không lên nhún vai: ``Chị ăn sạch trên đường về rồi còn đâu!''

Ryuuhou lập tức phồng mang trợn mắt: ``Phản bội! Chị đi làm phụ kho vật liệu mà ăn một mình không chia cho em gì cả!''

``Chị phải bỏ lao động chân tay với đếm từng thanh sắt ra mới đổi được đấy nhé! Người ta gọi đó là thành quả lao động chính đáng!'' tôi cãi lại đầy lý lẽ.

Kaigou từ trong bếp bước ra, nghe vậy liền gật gù tán đồng: ``Lý lẽ nghe rất vững đấy.''

``Kaigou-nee!! Chị ấy ăn một mình mà chị còn bênh nữa!'' Ryuuhou kêu gào phản đối.

Cô em mười lăm tuổi bèn ấm ức đi lại sofa, tính bật TV xem thể thao cho khuây khỏa. Nhưng vừa sờ tay xuống khe ghế, Ryuuhou đã reo lên: ``Ơ! Cái remote TV sao lại nằm dưới kẽ đệm ghế thế này?! Kson-san, lại là chị nhét vào đúng không?!''

Tôi giật nảy người, lập tức nhảy dựng lên phân trần: ``Chị nhét hồi nào! Lúc sáng chị đi phụ kho nó đã ở đó rồi mà! Mới lại đưa đây, chị đang định xem tiếp chương trình giải trí!''

``Không đời nào! Em tìm thấy trước thì em xem trước! Trận bóng đá derby sắp bắt đầu rồi, chị nhường em đi!'' Ryuuhou ôm khăng khăng chiếc điều khiển vào lòng.

Tôi rón rén lao tới định chộp lấy: ``Chị đi lao động khuân vác vất vả cả ngày về, phải được ưu tiên giải trí thư giãn chứ! Đưa đây cho chị!''

``\textbf{Phụ kho vật liệu có một buổi sáng mà chị làm như đi vận hành công trường xây dựng không bằng ấy! Trả remote cho em!!}'' Ryuuhou gào lên, hai chị em bắt đầu giằng co chiếc remote qua lại ngay trên ghế sofa, tiếng gầm gừ chí chố vang khắp phòng khách.

*XOẠT!*

Đúng lúc cuộc chiến đang ở hồi gay cấn nhất, một bàn tay dứt khoát vươn ra. Kaigou xuất hiện ngay bên cạnh từ bao giờ. Không một động thái thừa, cô tung một chiêu rút phắt chiếc remote khỏi tay cả tôi lẫn Ryuuhou chỉ trong một vệt chuyển động nhanh như chớp.

Chị Tổng giơ cao chiếc remote lên không trung, ánh mắt sắc lẹm lướt qua hai kẻ vừa tranh giành: ``Remote bị tịch thu cho đến khi cả hai người học được cách chia sẻ một cái TV mà không gào lên như ở ngoài bãi vật liệu. Ai có ý kiến gì không?''

Ryuuhou và tôi cùng há hốc miệng, đồng thanh gào lên đầy uất hận: ``\textbf{Kaigou-nee / Kaigou-paisen! Thế là bất công cực kỳ luôn đấy!!}''

Kaigou thản nhiên nhét tọt chiếc remote vào túi tạp dề, quay lưng đi thẳng vào bếp: ``Khiếu nại không có hiệu lực. Hai người có thể chọn ngồi nói chuyện hòa bình với nhau, hoặc ra sân sau vác bao xi măng chạy năm vòng cho hạ nhiệt rồi tính tiếp.''

Delutaya ngồi bên cạnh chứng kiến toàn bộ màn phán xử nhanh gọn lẹ liền bật cười rúc rích. Hai chúng tôi nhìn nhau, rồi cùng thở dài thượt ngả ngửa ra lưng ghế sofa.

Tôi tháo chiếc thẻ nhân viên tạm thời đeo trên cổ áo ra, nhìn dòng chữ \textbf{KSON - PHỤ KHO} được viết bằng bút đánh dấu cẩn thận trên tấm thẻ nhựa bảo hộ.

Một tấm thẻ tạm, vài giờ lao động điều phối tôn thép, một món quà ấm áp từ công trình và một ngôi nhà đầy ồn ào nhưng ngập tràn tiếng cười. Ngày hôm nay chẳng cần phải trở thành điều gì đó quá đỗi phi thường. Nó chỉ đơn giản là một ngày sống tử tế, có ích và trọn vẹn, và như thế là quá đủ rồi.

\end{document}